% TOP_PROBLEM: {"id": 30006687, "title": "P versus NP ∩ coNP", "category_id": 15, "category": {"id": 15, "name": "computer_science", "display_name": "Computer Science", "description": "Computational complexity, algorithms, and theoretical CS.", "slug": "computer-science", "order_index": 15, "created_at": "2026-07-31T15:26:25.671Z"}, "status": "open"}
% FIRST_POSTED: 2026-09-27
% !TeX program = lualatex
% Standalone notebook; compile twice: lualatex 30006687.tex
% Original section preserved verbatim except for marked insertion blocks.
\documentclass[11pt,a4paper]{article}
\usepackage[margin=24mm,headheight=14pt]{geometry}
\usepackage{fontspec}
\setmainfont{Latin Modern Roman}
\setsansfont{Latin Modern Sans}
\setmonofont{Latin Modern Mono}
\usepackage{amsmath,amssymb,mathtools}
\usepackage{unicode-math}
\setmathfont{Latin Modern Math}
\usepackage{microtype}
\usepackage{enumitem,xurl,needspace}
\usepackage[dvipsnames]{xcolor}
\usepackage{hyperref}
\hypersetup{unicode=true,colorlinks=true,linkcolor=MidnightBlue,urlcolor=MidnightBlue,
 pdftitle={Research notebook: Problem 30006687},
 pdfauthor={Open Problems Math research notebook},bookmarksnumbered=true}
\usepackage{bookmark,titlesec,fancyhdr}
\titleformat{\section}{\Large\bfseries\raggedright}{\thesection}{0.7em}{}
\pagestyle{fancy}\fancyhf{}
\fancyhead[L]{\small\sffamily Open Problems Math\enspace|\enspace Research notebook}
\fancyhead[R]{\small\sffamily Problem 30006687}
\fancyfoot[L]{\small\sffamily 27 September 2026}
\fancyfoot[R]{\small\sffamily \thepage}
\setlength{\parindent}{0pt}
\setlength{\parskip}{5pt plus 1pt}
\setlength{\emergencystretch}{3em}
\setcounter{secnumdepth}{1}
\setcounter{tocdepth}{1}
\setlist[itemize]{leftmargin=3em,itemsep=5pt,topsep=4pt,parsep=0pt}
\allowdisplaybreaks
\newcommand{\N}{\mathbb{N}}
\newcommand{\Z}{\mathbb{Z}}
\newcommand{\Q}{\mathbb{Q}}
\newcommand{\R}{\mathbb{R}}
\newcommand{\C}{\mathbb{C}}
\newcommand{\F}{\mathbb{F}}
\newcommand{\A}{\mathbb{A}}
\newcommand{\PP}{\mathbb P}
\newcommand{\E}{\mathbb{E}}
\newcommand{\eps}{\varepsilon}
\newcommand{\dd}{\,\mathrm d}
\newcommand{\sref}[2]{\hyperlink{source:#1:#2}{[S#2]}}
\newcommand{\eref}[2]{\hyperlink{extra:#1:#2}{[E#2]}}
\newif\ifshowcatalogue
\showcataloguetrue
\newcommand{\cataloguescope}[1]{\ifshowcatalogue
 \paragraph{Exact scope in the supplied catalogue.}
 \begin{quote}\small #1\end{quote}\fi}
\numberwithin{equation}{section}
\begin{document}
\setcounter{section}{0}
% BEGIN PRESERVED SECTION
% ================================================================
% problemId: 30006687
% Canonical problem ID: 30006687
% exactTarget SHA-256: 08f47bb409d1291266968f990686fc8b612ecc5632fbbc5bc0e36003c6c487cc
\clearpage
\section*{\texorpdfstring{P versus NP ∩ coNP}{P versus NP ∩ coNP}}\label{30006687}
\subsection{Definitions and mathematical statement}
A language lies in $\mathsf{NP}\cap\mathsf{coNP}$ when both membership and nonmembership have polynomial-length certificates verifiable in deterministic polynomial time. Thus there are verifiers $V_1,V_0$ and polynomial length bounds satisfying
\[
 x\in L\iff\exists y\,V_1(x,y)=1,\qquad
 x\notin L\iff\exists z\,V_0(x,z)=1.
\]
Determine whether every such language belongs to $\mathsf P$, that is, $\mathsf P=\mathsf{NP}\cap\mathsf{coNP}$. Algorithms are uniform and unrelativized, and inputs are total languages rather than promise problems.
\par\smallskip\textit{Formulation and concept references:} \sref{30006687}{1}.
\cataloguescope{Determine whether P = NP \ensuremath{\cap} coNP for ordinary uniform deterministic and nondeterministic polynomial-time Turing-machine classes of total languages over a finite alphabet; equivalently, whether every language with polynomially checkable certificates for both membership and nonmembership has a deterministic polynomial-time decision algorithm.}
\subsection{Short English statement}
When both answers have efficiently checkable certificates, must the answer itself be efficiently computable?
% BEGIN RESEARCH 34
\subsection{Research attempt}

\paragraph{Current frontier.}
The language-class question is exactly the dual-certification question in
Aaronson's lecture, not a promise or search-class comparison.
\sref{30006687}{1}. Hitchcock--Sekoni--Shafei obtain the corresponding separation
relative to individually random permutation oracles; their Theorem~5.1 does
not construct an ordinary polynomial-time permutation with a hard inverse.
\sref{30006687}{3}, \eref{30006687}{1}. Recent work on total randomized search separates
several black-box models and gives conditional white-box consequences, but
it does not turn totality into deterministic polynomial-time search.
\eref{30006687}{2}. The literature points to two distinct obstacles: certificates
may be difficult to locate, and an oracle carrying the desired hard data is
not an effective construction of those data. The recent theorem statements
were inspected, not independently proof-audited in full.

\paragraph{Attack route.}
For a negative answer, a natural target is a uniquely determined object whose
bits are efficiently certifiable in both directions, such as the inverse of
a polynomial-time permutation. For a positive answer, one could try to
combine the two certificate families into an efficient search. I pursue the
common interface between these routes: unique \emph{outputs} versus unique
\emph{certificates}, and the complexity of selecting a certificate consistent
with partial information. These distinctions are essential to avoid assuming
the answer in a canonicalization step.

\paragraph{Attempt.}
Call a relation $R(x,y,z)$ a \emph{total single-output verification scheme}
if it is deterministic polynomial-time decidable, the lengths of $y,z$ are
bounded by fixed polynomials in $|x|$, and for every $x$ there is exactly one
string $y=f(x)$ for which some $z$ satisfies $R(x,y,z)$. Pad the output to a
fixed polynomial length, with an encoded original length if necessary.
There may be many valid $z$ for this one $y$.

\textbf{Lemma 34.1 (an exact output-search formulation).}
The equality $\mathsf P=\mathsf{NP}\cap\mathsf{coNP}$ holds if and only if
every total single-output verification scheme has its output function in
$\mathsf{FP}$.

\emph{Proof.} Given such a scheme, let
\[
 B=\{\langle x,i\rangle:1\le i\le\ell(x),\ f(x)_i=1\}.
\]
On a valid input, a certificate for membership guesses $y,z$ with
$R(x,y,z)$ and $y_i=1$; a certificate for nonmembership guesses a valid
output with $y_i=0$. Uniqueness and totality prove both soundness and
completeness. Malformed inputs are rejected deterministically and so also
have a polynomial-time complementary verifier. Thus
$B\in\mathsf{NP}\cap\mathsf{coNP}$. Under the stated equality, deciding
its polynomially many bits computes $f$ in polynomial time.

Conversely, for $L\in\mathsf{NP}\cap\mathsf{coNP}$ with verifiers $V_1,V_0$,
let $R(x,b,z)=V_b(x,z)$ for the one-bit output $b\in\{0,1\}$. Exactly one
value of $b$ has a certificate. An output algorithm for this scheme decides
$L$. This proves the equivalence with all quantifiers over fixed uniform
verifiers and fixed polynomial bounds intact.\hfill$\square$

The lemma identifies the requisite search power, but does not establish it.
For an arbitrary total polynomial-time relation, independently choosing each
bit that appears in \emph{some} solution is invalid. With solution set
$\{01,10\}$ the existential bit answers are $11$, not a solution. Sequential
prefix extension repairs this consistency problem, but introduces the
predicate
\begin{equation}\label{r34:prefix}
 E(x,u)\quad\Longleftrightarrow\quad
 \exists y,z\,[R(x,y,z)\ \text{and }u\text{ is a prefix of }y].
\end{equation}
For multiple possible outputs, nonextendibility does not in general have a
certificate obtained just by displaying another solution. For a unique output,
a conflicting output does certify nonextendibility, so $E$ is again in
$\mathsf{NP}\cap\mathsf{coNP}$. This recovers the same unresolved decision
problem rather than eliminating it.

\medskip
\textbf{Trying to force a canonical certificate.}
One might replace a verifier $V(x,w)$ by its lexicographically least valid
witness. Correct verification would then require
\[
 V(x,w)=1\quad\text{and}\quad
              \forall w'<w,\ V(x,w')=0.
\]
The second condition is universal. A polynomial-time verification algorithm
for it has not been furnished merely by the existence of dual certificates
for $x$. Even when $x$ is known to be a yes-instance, its no-verifier need
not certify that a \emph{subcollection} of yes-certificates is empty.
The following elementary decision-tree comparison makes the algorithmic
cost of this missing operation explicit.

\textbf{Lemma 34.2 (what a certificate-elimination argument actually uses).}
Let $f:\{0,1\}^n\to\{0,1\}$ be total. A bit certificate is a partial
assignment forcing an output. Suppose the largest minimum zero-certificate
size is $c_0$ and the analogous one-size is $c_1$. Then a deterministic
decision tree of depth at most $c_0c_1$ exists. This bound does not charge for
finding the certificates defining the tree.

\emph{Proof.} If $f$ is constant, no queries are needed. Otherwise choose
a one-input and a one-certificate for it of size at most $c_1$, and query
those positions. If every value matches, output one. On each other branch,
consider a zero-input still consistent with the answers and one of its
zero-certificates of size at most $c_0$. That certificate must conflict with
the chosen one-certificate: otherwise a common completion would force both
outputs. At least one of its positions has now been queried. Deleting the
queried positions leaves a zero-certificate of size at most $c_0-1$ for the
restricted function. Its one-certificate complexity is at most $c_1$.
Induction proves the depth bound; if one output type disappears, the branch
is constant.\hfill$\square$

This attractive bound does not solve the Turing-machine problem. Selecting
a consistent one-input can be difficult even when evaluating $f$ is easy.
Given a Boolean circuit $C$ and a partial assignment $u$, the task of finding
any $a$ extending $u$ with $C(a)=1$, or correctly reporting none, decides
CircuitSAT by taking $u$ empty. In contrast, evaluating $C(a)$ on a given
$a$ is polynomial time. Thus the decision-tree construction can hide an
NP-hard \emph{selection} subroutine inside an argument about an easy decision
function. Moreover arbitrary NP certificates are not necessarily partial
assignments of input bits. Both mismatches must be repaired, not suppressed.
The CircuitSAT formulation is given in \sref{30006687}{1}.

\medskip
\textbf{Testing the negative route.}
Let $p_n$ be any uniformly polynomial-time bijection of $\{0,1\}^n$.
The language
\[
 L_p=\{\langle1^n,y,i\rangle:(p_n^{-1}(y))_i=1\}
\]
uses the inverse $x$ as a certificate for either bit. It lies even in
$\mathsf{UP}\cap\mathsf{coUP}$, because $x$ is unique. If $L_p$ has a
deterministic polynomial-time algorithm, then $n$ bit queries invert $p_n$
on every input. Conversely a polynomial-time inverse decides $L_p$.
Therefore an explicit polynomial-time permutation without a polynomial-time
inverse would separate the classes in this entry. Average-case one-wayness,
as in UnsolvedMath problem~30006696, is more than this deduction needs; worst-case inverse
hardness suffices. No such unconditional hardness is proved here. The
oracle construction in \eref{30006687}{1} cannot replace it, since its forward
permutation is supplied as an oracle rather than a uniform algorithm.

\paragraph{Stress test.}
Lemma~34.1 permits many auxiliary certificates for one output, so it does not
silently replace $\mathsf{NP}\cap\mathsf{coNP}$ by its unambiguous subclass.
Uniqueness concerns the entire output, not merely each separately selected
bit. Totality is indispensable to the complementary bit verifier; for a
partial function, an outputless input has neither bit certificate. Constant
functions and malformed encodings have been treated explicitly.

The accompanying script checks the inequality $D(f)\le C_0(f)C_1(f)$ for
all Boolean functions of three variables, computing optimal decision-tree
depth by recursion on partial assignments. It also checks the two-witness
bit-selection failure. These finite checks say nothing about the uniform
cost of constructing trees as the input length grows. A proof of separation
here would imply $\mathsf P\ne\mathsf{NP}$; the exact-output equivalence
cannot be used to obtain that implication without supplying the missing
hard function.

\Needspace{8\baselineskip}
\paragraph{Outcome.}
\textbf{Outcome: Exploratory approach with unresolved gap.}

The attempt isolates an exact single-output search reformulation and verifies
why two tempting positive strategies fail: independent witness-bit choice
can be inconsistent, and canonical or decision-tree certificate selection
can conceal a harder problem than the original decision task. The negative
route is reduced to an explicit worst-case-hard permutation inverse, but
no such construction is established. These are proved structural deductions,
not a new class separation or a claim of historical novelty.

% END RESEARCH 34

\subsection{Sources}
\begin{itemize}\raggedright
\item[\textnormal{[S1]}]\hypertarget{source:30006687:1}{} Scott Aaronson, PHYS771 Lecture 6: P, NP, and Friends.
\url{https://www.scottaaronson.com/democritus/lec6.html}.
{\small\textit{Catalogue formulation source.}}
\item[\textnormal{[S2]}]\hypertarget{source:30006687:2}{} Very few problems are in NP intersect coNP but not known to be in P. What to make of that?.
\url{https://blog.computationalcomplexity.org/2024/09/very-few-problems-are-in-np-intersect.html?m=0}.
{\small\textit{Catalogue research/status reference; not a proof endorsement.}}
\item[\textnormal{[S3]}]\hypertarget{source:30006687:3}{} Random Permutations in Computational Complexity.
\url{https://arxiv.org/abs/2511.08786}.
{\small\textit{Catalogue research/status reference; not a proof endorsement.}}
% BEGIN ADDED REFERENCES 34
\item[\textnormal{[E1]}]\hypertarget{extra:30006687:1}{} John M. Hitchcock,
Adewale Sekoni and Hadi Shafei, \emph{Random Permutations in Computational
Complexity}, arXiv:2511.08786v1 (2025), abstract and Theorem~5.1.
\url{https://arxiv.org/abs/2511.08786v1}.
{\small\textit{Version-specific supplement to [S3]; oracle-relative theorem,
checked 27 September 2026.}}
\item[\textnormal{[E2]}]\hypertarget{extra:30006687:2}{} Noah Fleming, Stefan
Grosser, Siddhartha Jain, Jiawei Li, Hanlin Ren, Morgan Shirley and
Weiqiang Yuan, \emph{Total Search Problems in $\mathsf{ZPP}$},
arXiv:2512.01138 (2025), abstract and scope of the total-search models.
\url{https://arxiv.org/abs/2512.01138}.
{\small\textit{Statement-level frontier comparison, checked
27 September 2026; not used as a premise of the notebook lemmas.}}

% END ADDED REFERENCES 34
\end{itemize}

% END PRESERVED SECTION
\end{document}
